\documentclass[11pt,letterpaper,twocolumn]{article}
\usepackage[utf8]{inputenc}
 \renewcommand{\familydefault}{\sfdefault}
\usepackage[spanish]{babel}
\usepackage{amsmath}   		
\usepackage{amsfonts}  		
\usepackage{amssymb}
\usepackage{graphicx}
\usepackage{enumitem}
\usepackage[bookmarks=true,colorlinks=true,linkcolor=black,citecolor=black,menucolor=black,urlcolor=black]{hyperref}  
\frenchspacing
\addto\captionsspanish{\renewcommand*{\tablename}{\textbf{Tabla}}} 
\numberwithin{equation}{section}
\numberwithin{figure}{section} 
\textwidth=17cm
\textheight=21cm
\topmargin=-1cm
\oddsidemargin=0cm
%\setlength{\columnsep}{0.7cm}	% la separaion de las columnas
\newcommand{\nl}{~\\}
\newcommand{\ohm}{\Omega}
\begin{document}
% Logo, Resumen y weas
%\begin{algo}
% Logo, titulo, y cuestiones asi
\twocolumn[
\begin{@twocolumnfalse}
% Logo y cuestiones
\begin{tabular}{cl}
\begin{tabular}{l}
%{ \includegraphics[scale=0.35]{logo-uc-color.jpeg} }
{ \includegraphics[scale=0.28]{logo-uc-byn.png} }
\end{tabular}
& 
\begin{tabular}{l}
\nl
Pontificia Universidad Católica de Chile \nl
Facultad de Física - Departamento de Física  \nl
FIZ0421 - F\'{i}sica Experimental Avanzada \nl
Profesor: 
\end{tabular}
\end{tabular}

\vspace{0.1cm}

% Ahora el nombre del experimento, los nombres y la fecha
\begin{center}
{\textbf{\huge{Experimento X}}} 
\\
\vspace{0.2cm}
Autores 
\vspace{0.1cm}

%Fecha de entrega
Fecha de Entrega
\end{center}
% aca va el resumen
\begin{abstract}
{\normalsize
Se estudio un plasma ...
} \nl \nl
\end{abstract}		
\end{@twocolumnfalse}]
%\end{algo}

\section{Introducci\'on}
%Aqui va la teoria y las formulas 
\section{Montaje Experimental}
%Dibujos y descripcion de aparataje
\section{Procedimiento}
%Lo que se hizo y como se usaron las formulas
\section{Resultados}
%Los resultados, graficos y tablas
\section{An\'alisis}
%Por que esos resultados y sus errores
\section{Conclusiones}
%Evaluar si se cumplio el objetivo
%Puntos debiles y fuertes (errores por ej)
%Posibles mejoras
\begin{thebibliography}{99}

\end{thebibliography}
% La bibliografia
%\newpage
%\begin{thebibliography}{99}
%\bibitem{ref:LibroChicoLab} David N. Ruzie. (1994). Electric Probes For Low Temperatures Plasmas. New York, New York: The American Vacuum Society.
%\\bibitem{ref:PlasmaEngineering} Reece Roth. (1995). Industrial Plasma Engineering. London: The Institute of Physics.
%\end{thebibliography}

\end{document}

